%%%%%%
%
% $Autor:  Wings $
% $Datum: 25/02/04 $
% $Directory: Contents\en\DVC.tex $
% $Version: 2 $
% $Autor:  $
% $Review date: 25/02/04 $
%
%%%%%%

\chapter{Data Version Control}

	
\section{Introduction}

Data Version Control (DVC) has been utilized to efficiently manage and track changes in the dataset throughout the development process \cite{Iterative:2020}.
	
\section{Plan}

The project workflow includes the following stages:

\begin{enumerate}
  \item Project Initialization:
    \begin{itemize}
	  \item Set Up Version Control:
	  \item Initialize a Git repository for your project.
	  \item Install DVC and initialize it within the repository.
%	  \item Configure a remote storage (e.g., AWS S3, Google Drive) to store large datasets and model files.
	\end{itemize}

  \item Dataset Management:
	\begin{itemize}
		\item Use DVC to add the dataset to version control, enabling tracking of changes and facilitating dataset versioning.
		\item Commit the DVC files and dataset changes to the Git repository.
	\end{itemize}

  \item Data Preprocessing:
	\begin{itemize}
		\item 	Define a DVC pipeline that includes stages for data preprocessing, feature extraction, and any other necessary steps.
		\item This approach ensures that any changes in the preprocessing steps are tracked and can be reproduced.
	\end{itemize}

	
  \item Model Training:
    \begin{itemize}
    	\item Create a DVC pipeline stage for model training, specifying dependencies such as the preprocessed data and model architecture.
    	\item This setup allows for easy reproduction of training experiments and comparison of different model versions.
    \end{itemize}

	
	\item Model Evaluation:
	\begin{itemize}
		\item Utilize DVC to track different training experiments, including hyperparameter variations and model architectures.
		\item This practice facilitates the identification of the most effective model configurations
	\end{itemize}

	
	\item Deployment Preparation:
	\begin{itemize}
		\item Use DVC to version the optimized model, ensuring that the exact model version used in deployment is tracked and reproducible.
	\end{itemize}

	
	\item Documentation and Collaboration:

	DVC's tracking capabilities assist in maintaining a clear history of changes, which is crucial for collaboration and future reference.
\end{enumerate}
	
\section{Installation}

To replicate the project setup, follow the steps below.
	
	
\subsection{Detailed Installation and usage Steps}

	\begin{itemize}
		
		\item Project Initialization:
		Initialize a Git repository to manage your project's source code.
		
\medskip

\SHELL{git init MyProject}

\SHELL{cd MyProject}

\medskip
		
Install required Python libraries:

\SHELL{pip install -r requirements.txt}

\medskip

Install DVC:

\SHELL{pip install dvc}

\medskip
		
\begin{center}
  \includegraphics[width=0.8\textwidth]{DVC/DVCInstallation.png}
  \captionof{figure}{Data Version Control Installation in the Command line}
			\label{fig:DVC_Installation}
\end{center}

\end{itemize}

\medskip
		
\subsection{Initialize DVC}

Set up DVC within your repository to manage data and model versioning.

\medskip


\SHELL{dvc init}

\medskip


\begin{itemize}
	\item Dataset Management: Acquire the Google Speech Commands dataset, which contains 65,000 one-second audio files of 30 short words, contributed by thousands of individuals.
		\item Organize the Data: Place the dataset in a directory named data/.
		
		\item Track the Dataset with DVC: Use DVC to monitor changes in the dataset.
		
		\medskip
		
		\SHELL{dvc add data/}
	
		\item Commit Changes: Record the dataset tracking in Git.
		
		 \medskip
		 
		 
		\SHELL{git add data.dvc .gitignore}
		
		\SHELL{git commit -m "{}Add Google Speech Commands dataset with DVC tracking"}

        \medskip
		
		\item Data Processing Pipeline : Create a DVC pipeline to preprocess the audio data, such as converting audio formats, normalizing, and feature extraction (e.g., Mel-Frequency Cepstral Coefficients).
		
		\item Create a DVC Pipeline File: Define the stages in a dvc.yaml file.
		
	\begin{lstlisting}[language=python]
		stages:
		preprocess:
		cmd: python scripts/preprocess.py
		deps:
		- scripts/preprocess.py
		- data/
		outs:
		- data/processed/
	\end{lstlisting}
	
	\item Commit Pipeline Definition: Save the pipeline configuration.
			
	\medskip
	
	
	\SHELL{git add dvc.yaml}
	
	\SHELL{git commit -m "Define data preprocessing pipeline"}
	
	\medskip
	
	
	\item Model Training Pipeline: Add a stage in the DVC pipeline for model training.
	
	\begin{lstlisting}[language=python]
		stages:
		train:
		cmd: python scripts/train.py
		deps:
		- scripts/train.py
		- data/processed/
		outs:
		- models/kws_model.pth
	\end{lstlisting}
	
	\item Commit Training Stage: Update the pipeline configuration.
			
	\medskip
	
	
	\SHELL{git add dvc.yaml}
	
	\SHELL{git commit -m "Add model training stage to pipeline"}
	
	\medskip
	
	
	\item Model Evaluation and Deployment: Add a stage to evaluate the trained model.
	
	\begin{lstlisting}[language=python]
		stages:
		evaluate:
		cmd: python scripts/evaluate.py
		deps:
		- scripts/evaluate.py
		- models/kws_model.pth
		outs:
		- reports/metrics.json
	\end{lstlisting}
	
	\item Commit Evaluation Stage: Save the evaluation configuration.
			
	\medskip
	
	
	\SHELL{git add dvc.yaml}
	
	\SHELL{git commit -m "Add model evaluation stage to pipeline"}
	
	\medskip
	
	
	\item Set Up Remote Storage: Configure DVC to use remote storage (e.g., AWS S3, Google Drive) for storing large files and datasets.
			
	\medskip
	
	
	\SHELL{dvc remote add -d myremote s3://mybucket/dvcstore}
	
	\medskip
	
	
	\item Push Data to Remote: Upload tracked data and models to the remote storage.
			
	\medskip
	
	
	\SHELL{dvc push}
	
	\medskip
	
	
	\item Share Project with Collaborators: Ensure that collaborators can reproduce the environment and results by pulling data from the remote storage.
			
	\medskip
	
	
	\SHELL{git clone https://github.com/yourusername/MyProject.git}
	
	\SHELL{cd MyProject}
	
	\SHELL{dvc pull}
	
	\medskip
	
	
	\item Track Dataset Changes: Use DVC to monitor and manage any updates or changes to the dataset.
	
	\item Update Pipelines as Needed: Modify the DVC pipelines to accommodate new preprocessing steps, model architectures, or evaluation metrics.
	
	\end{itemize}
	
	The study \href{https://ieeexplore.ieee.org/document/9425888}{On the Co-evolution of ML Pipelines and Source Code - Empirical Study of DVC Projects} examines the integration of DVC in machine learning projects and highlights the importance of versioning both data and code to maintain consistency and reproducibility.
	\Mynote{cite!}
	
	By following this plan, you can effectively manage and track changes in your dataset throughout the development of your project, ensuring a reproducible and collaborative workflow.

	
	
